\chapter{Propriétés angulaires des parallèles}


\begin{itemize}
\it Démonter que pour que deux droites soient sécantes, il suffit : 
\begin{enumerate}
\item Que l'une soit perpendiculaire et l'autre oblique par rapport à une troisième.
\item Qu'elles soient perpendiculaires aux côtés d'un angle saillant.
\end{enumerate}
\it Étant données deux parallèles coupées par une sécante, montrer que : 
\begin{enumerate}
\item Les bissectrices de deux angles alternes-internes ou correspondants sont parallèles. 
\item Les bissectrices de deux angles intérieurs d'un même côté sont perpendiculaires. 
\item Énoncer les réciproques de ces propriétés et les démontrer. 
\end{enumerate}
\it Sur les côtés d'un angle de sommet $O$ on porte deux longueurs égales $OA=OB$. Puis extérieurement à cet angle, on construit deux angles égaux $\widehat{OAx}$ et $\widehat{OBy}$ et on porte sur $[Ax)$ et $[By)$ deux longueurs égales $AC=BD$.\begin{enumerate}
\item Comparer les triangles $OAC$ et $OBD$.
\item Montrer que les angles $\widehat{AOB}$ et $\widehat{COD}$ ont même bissectrice.
\item Comparer les directions de $(AB)$ et $(CD)$.
\end{enumerate}
\it On considère deux angles adjacents supplémentaires $\widehat{AOB}$ et $\widehat{BOC}$ et leurs bissectrices $[Ox)$ et $[Oy)$. Par le point $B$, on mène la parallèle à $(AC)$ qui coupe $[Ox)$ en $M$ et $[Oy)$ en $N$.\begin{enumerate}
\item Comparer $MB$ et $OB$ ; puis $NB$ et $OB$.
\item Que représente le point $B$ pour le segment $[MN]$ ?
\end{enumerate}
\it Par le point de rencontre $I$ des bissectrices intérieures des angles en $B$ et $C$ du triangle $ABC$, on mène la parallèle à $(BC)$, qui coupe $(AB)$ en $D$ et $(AC)$ en $E$. \begin{enumerate}
\item Comparer $DB$ et $DI$ puis $EC$ et $EI$.
\item Montrer que $DE=BD+CE$. 
\end{enumerate}
\it Soit un triangle $ABC$ rectangle en $A$. Sur la perpendiculaire en $C$ à $(AC)$ on porte des segments $[CD]$ et $[CE]$ égaux à $[BC]$.\begin{enumerate}
\item Comparer les directions de $(AB)$ et $(DE)$.
\item Que représentent $(BD)$ et $(BE)$ pour l'angle en $B$ ? 
\end{enumerate}
\it Soit un triangle $ABC$. On mène les bissectrices intérieures des angles en $B$ et $C$ qui coupent en $D$ et $E$ la parallèle menée par $A$ à $(BC)$. \begin{enumerate}
\item Comparer $AD$ et $AB$ puis $AE$ et $AC$.
\item Montrer que $DE=AB+AC$.
\item Reprendre le problème avec les bissectrices extérieures.
\end{enumerate}
\it Soit un triangle $ABC$. On mène par le milieu $D$ de $[BC]$ la perpendiculaire à la bissectrice intérieure de l'angle en $A$. Cette perpendiculaire coupe $(AB)$ en $E$ et $(AC)$ en $F$. $(EF)$ coupe alors la parallèle menée par $B$ à $(AC)$ en $G$. 
\begin{enumerate}
\item Montrer que les triangles $AEF$ et $BEG$ sont isocèles.
\item Comparer les triangles $DBG$ et $DCF$. 
\item Démontrer l'égalité des longueurs $BE$ et $CF$.
\end{enumerate}



\end{itemize}